\documentclass{article}
\usepackage[T1]{fontenc}
\usepackage{lmodern}
\usepackage{graphicx}
\usepackage{amsmath,amssymb}
\usepackage[a4paper,margin=2cm]{geometry}
\usepackage[absolute,overlay]{textpos}
\setlength{\TPHorizModule}{1mm}
\setlength{\TPVertModule}{1mm}
\usepackage[indonesian]{babel}
\usepackage{hyperref}
\setlength{\emergencystretch}{2em}
% unit: Halaman 19

\begin{document}

% === Halaman 19 (python/native) ===

Dekripsi

• Dekripsi dilakukan dengan menggunakan kunci privat. 

• Mula-mula penerima pesan menghitung n–1 , yaitu balikan dari n 

modulo m, sedemikian sehingga

   n  n–1  1 (mod m).  

• Kalikan setiap kriptogram dengan n–1 , lalu nyatakan hasil kalinya 

sebagai penjumlahan elemen-elemen kunci privat  untuk 

memperoleh plainteks  dengan menggunakan algoritma pencarian 

solusi superincreasing knapsack.

19

\end{document}
